\documentclass[11pt,a4paper]{article}
\usepackage[utf8]{inputenc}
\usepackage[T1]{fontenc}
\usepackage{lmodern}
\usepackage[french]{babel}
\usepackage{amsmath,amssymb,mathtools}
\usepackage{icomma}
\usepackage{xcolor}
\usepackage{tikz}
\usepackage{pgfplots}
\usepackage[most]{tcolorbox}
\usepackage{enumitem}
\usepackage{fancyhdr}
\usepackage[hidelinks]{hyperref}
\usepackage{titlesec}
\usepackage[a4paper,margin=2.2cm,headheight=26pt,headsep=10pt]{geometry}
\usetikzlibrary{arrows.meta,positioning,calc,intersections,angles,quotes}
\usepgfplotslibrary{fillbetween}
\pgfplotsset{compat=1.18}
\definecolor{primary}{RGB}{30,75,145}
\definecolor{primarydark}{RGB}{20,50,100}
\definecolor{defcol}{RGB}{0,114,178}
\definecolor{thmcol}{RGB}{0,135,110}
\definecolor{excol}{RGB}{0,130,85}
\definecolor{remarkcol}{RGB}{120,70,180}
\definecolor{attncol}{RGB}{200,40,55}
\definecolor{methcol}{RGB}{85,85,100}
\definecolor{areacol}{RGB}{120,160,230}
\definecolor{areacol2}{RGB}{230,150,80}
\newcommand{\dd}{\mathrm{d}}
\newcommand{\R}{\mathbb{R}}
\newcommand{\ua}{\,\text{u.a.}}
\newcommand{\tg}{\tan}\newcommand{\cotg}{\cot}
\tcbset{boxbase/.style={enhanced,breakable,boxrule=0.9pt,arc=2.5pt,left=3mm,right=3mm,top=2.3mm,bottom=2.3mm,fonttitle=\bfseries,coltitle=white,toptitle=0.6mm,bottomtitle=0.6mm}}
\newtcolorbox{methode}[1][]{boxbase,colback=methcol!7,colframe=methcol,sharp corners,title={\textsc{Comment faire}\ifx\relax#1\relax\relax\else\textmd{ : \itshape #1}\fi}}
\newtcolorbox{exemple}[1][]{boxbase,colback=excol!5,colframe=excol,title={\textsc{Exemple}\ifx\relax#1\relax\relax\else\textmd{ : \itshape #1}\fi}}
\newtcolorbox{idee}[1][]{boxbase,colback=defcol!5,colframe=defcol,title={\textsc{Idée clé}\ifx\relax#1\relax\relax\else\textmd{ : \itshape #1}\fi}}
\newtcolorbox{remarque}[1][]{boxbase,colback=remarkcol!6,colframe=remarkcol,title={\textsc{Remarque}\ifx\relax#1\relax\relax\else\textmd{ : \itshape #1}\fi}}
\newtcolorbox{attention}[1][]{boxbase,colback=attncol!6,colframe=attncol,title={\textsc{Attention}\ifx\relax#1\relax\relax\else\textmd{ : \itshape #1}\fi}}
\newtcolorbox{exo}[1][]{boxbase,colback={primary!4},colframe=primary,title={\textsc{Exercice #1}}}
\newtcolorbox{corr}[1][]{boxbase,colback={thmcol!5},colframe=thmcol,title={\textsc{Corrigé #1}}}
\tikzset{axe/.style={->,>=Stealth,black!65,line width=0.7pt},courbe/.style={primary,line width=1.3pt},courbe2/.style={areacol2!90!black,line width=1.3pt},dvert/.style={gray,dashed,line width=0.7pt}}
\pagestyle{fancy}\fancyhf{}
\fancyhead[L]{\small\textcolor{primarydark}{\textbf{Fiche 7} — Intégrale d'une fonction continue}}
\fancyhead[R]{\small\textcolor{primarydark}{\textsc{Thème 2 — Intégration par parties}}}
\fancyfoot[C]{\small\thepage}
\renewcommand{\headrulewidth}{0.8pt}
\titleformat{\section}{\large\bfseries\color{primarydark}}{\thesection}{0.6em}{}
\begin{document}
\section*{Tutoriel — L'intégration par parties (IPP)}

Certaines fonctions n'ont \emph{aucune} primitive dans le formulaire du Thème~1 :
$x\,e^x$, $x\cos x$, $\ln x$, $\arctan x$\dots\ Ce sont des \textbf{produits} (ou des fonctions
« isolées » comme $\ln x$) que l'on intègre grâce à une seule technique universelle :
l'\textbf{intégration par parties}. Elle n'est rien d'autre que la dérivée d'un produit lue à l'envers.

\begin{idee}[la formule et d'où elle vient]
On part de la dérivée d'un produit : $(UV)'=U'V+UV'$. En passant aux différentielles puis en
intégrant, $\displaystyle\int \dd(UV)=\int V\,\dd U+\int U\,\dd V$, c'est-à-dire $UV=\int V\,\dd U+\int U\,\dd V$.
En isolant la deuxième intégrale :
\[
\boxed{\;\int U\,\dd V \;=\; U\,V-\int V\,\dd U.\;}
\]
On choisit $U$ (qu'on va \textbf{dériver}) et $V'$ (qu'on va \textbf{intégrer}, avec $\dd V=V'\,\dd x$).
Le pari : que la nouvelle intégrale $\int V\,\dd U$ soit \emph{plus simple} que celle de départ.
\end{idee}

\subsection*{Comment choisir $U$ et $V'$ ?}

Tout le succès tient au bon choix. On prend pour $U$ (la fonction que l'on dérive) celle qui apparaît
le plus à \textbf{gauche} dans l'ordre de priorité suivant — retenu par le mnémo \textbf{« L.I.A.T.E. »} :
\[
\underbrace{\text{L}}_{\ln,\ \log}\ >\
\underbrace{\text{I}}_{\arcsin,\ \arctan}\ >\
\underbrace{\text{A}}_{x^n,\ \text{polynômes}}\ >\
\underbrace{\text{T}}_{\sin,\ \cos}\ >\
\underbrace{\text{E}}_{e^x,\ a^x}.
\]
\textbf{Pourquoi cet ordre ?} Un polynôme \emph{dérivé} perd un degré (et finit par disparaître) ; un
$\ln$ ou un $\arctan$ dérivé devient une fraction algébrique bien plus commode ; à l'inverse, $\sin$,
$\cos$, $e^x$ s'\emph{intègrent} sans se compliquer, on les met donc en $V'$.

\begin{methode}[dérouler une IPP en 4 gestes]
\begin{enumerate}[leftmargin=1.8em,topsep=2pt,itemsep=1pt]
  \item \textbf{Choisir} $U$ (règle L.I.A.T.E.) et $V'$ (tout le reste) ; calculer $U'$ (donc $\dd U=U'\,\dd x$) et $V=\int V'\,\dd x$.
  \item \textbf{Appliquer} $\displaystyle\int U\,\dd V=UV-\int V\,\dd U$.
  \item \textbf{Simplifier} la nouvelle intégrale (parfois il faut une \emph{deuxième} IPP, ou résoudre une équation en $I$).
  \item \textbf{Vérifier} en dérivant le résultat : on doit retomber sur l'intégrande.
\end{enumerate}
\end{methode}

\begin{exemple}[le cas d'école : $\displaystyle\int x\,e^x\,\dd x$]
\textbf{Choix.} $U=x$ (algébrique, à dériver) et $V'=e^x$ (exponentielle, à intégrer), donc $U'=1$ et $V=e^x$.
\[
\int x\,e^x\,\dd x=U V-\int V\,\dd U=x\,e^x-\int e^x\cdot 1\,\dd x=x\,e^x-e^x+C=\boxed{(x-1)\,e^x+C.}
\]
\textbf{Vérification.} $\bigl[(x-1)e^x\bigr]'=e^x+(x-1)e^x=x\,e^x$. \checkmark\ Le produit de départ est devenu
l'intégrale \emph{immédiate} $\int e^x\,\dd x$ : c'est tout l'intérêt d'avoir dérivé $x$ en $1$.
\end{exemple}

\begin{exemple}[$\ln x$, qui n'a pas de primitive « directe »]
On écrit $\ln x=\underbrace{1}_{V'}\cdot\underbrace{\ln x}_{U}$ : ici $U=\ln x$ (priorité L !) et $V'=1$, donc
$U'=\tfrac1x$ et $V=x$.
\[
\int \ln x\,\dd x=x\ln x-\int x\cdot\frac1x\,\dd x=x\ln x-\int 1\,\dd x=\boxed{x\ln x-x+C.}
\]
\textbf{Vérification.} $\bigl[x\ln x-x\bigr]'=\ln x+x\cdot\tfrac1x-1=\ln x$. \checkmark
\end{exemple}

\begin{remarque}[quand l'IPP « tourne en boucle »]
Pour $\displaystyle I=\int e^x\sin x\,\dd x$, aucune IPP ne simplifie vraiment : deux IPP successives
ramènent la même intégrale $I$. On obtient alors une \emph{équation} du type $I=(\ldots)-I$ que l'on
\textbf{résout} pour trouver $I$. Reconnaître cette situation évite de tourner en rond (voir Difficile).
\end{remarque}

\begin{attention}[deux pièges classiques]
\begin{itemize}[leftmargin=1.5em,topsep=1pt,itemsep=1pt]
  \item \textbf{Le signe $-$} de la formule ($-\int V\,\dd U$) est oublié une fois sur deux : écrivez-le tout de suite.
  \item \textbf{Le mauvais choix} $U\leftrightarrow V'$ complique l'intégrale au lieu de la simplifier. En cas de doute,
        essayez ; si $\int V\,\dd U$ est \emph{pire} que l'intégrale de départ, inversez les rôles.
\end{itemize}
\end{attention}
\end{document}
